% =====================================================================
% 源程序文档 LaTeX 模板
% 使用说明：将 {{变量名}} 替换为实际内容，将代码粘贴到 Verbatim 环境中
% 编译：xelatex -interaction=nonstopmode 文件名.tex （运行两遍）
%
% 关键点：
% - 10pt + \small + 窄边距 → 每页自然约 50 行
% - 代码连续排列，LaTeX 自然分页，不要手动插入分页符
% - 前30页 + 后30页 = 正好60页
% - \parindent = 0（代码文档不缩进）
% =====================================================================
\documentclass[a4paper,10pt]{article}
\usepackage[UTF8]{ctex}
\usepackage[top=2cm,bottom=2cm,left=1.5cm,right=1.5cm]{geometry}
\usepackage{fancyhdr}
\usepackage{lastpage}
\usepackage{hyperref}
\usepackage{fancyvrb}

\hypersetup{hidelinks}

\pagestyle{fancy}
\fancyhf{}
\fancyhead[C]{{{软件全称}}V{{版本号}}\quad 源程序}
\fancyfoot[C]{第 \thepage\ 页\quad 共 \pageref{LastPage} 页}
\renewcommand{\headrulewidth}{0.5pt}
\renewcommand{\footrulewidth}{0.5pt}
\fancypagestyle{plain}{\pagestyle{fancy}}

\setlength{\parindent}{0em}

\begin{document}

\begin{center}
  {\LARGE\bfseries {{软件全称}}\quad 源程序}

  \vspace{0.5em}
  {\large 著作权人：{{著作权人名称}}}

  \vspace{0.3em}
  软件名称：{{软件全称}}\qquad 版本号：V{{版本号}}\qquad 生成日期：{{生成日期}}

  \vspace{0.3em}
  源程序量：{{源程序量}}行\qquad 本文档含前30页 + 后30页
\end{center}

\vspace{0.3em}

\begin{center}\large\bfseries —— 前30页（第1--{{前30页末行}}行）——\end{center}

\begin{Verbatim}[fontsize=\small]
% === 在此粘贴前30页的代码 ===
% 每个文件开头用注释标记：
% // ========== path/to/file.js ==========
% 代码连续排列，不要手动分页
\end{Verbatim}

\begin{center}\large\bfseries —— 后30页（第{{后30页起始行}}--{{最末行}}行）——\end{center}

\begin{Verbatim}[fontsize=\small]
% === 在此粘贴后30页的代码 ===
\end{Verbatim}

\end{document}
